\section{静态节点池：打破循环依赖}
\label{sec:vma-pool}

回顾初始化顺序（第~\ref{ch:kmalloc}~章）：

\codefile{src/kernel/mm/vma.c}
\begin{lstlisting}[style=cstyle]
pmm_init();        /* 1. 物理内存 */
vmm_init();        /* 2. 页表 */
kmalloc_init();    /* 3. 内核堆 */
vma_init();        /* 4. VMA — 依赖 kmalloc */
\end{lstlisting}

VMA 初始化在 \texttt{kmalloc\_init()} \emph{之后}——所以后端\emph{可以}用 \texttt{kmalloc}
分配内部节点。但存在一个微妙问题：

\begin{whybox}
如果 \texttt{kmalloc} 的堆增长（\texttt{heap\_grow}）需要调 \texttt{vmm\_alloc\_pages}，
而 \texttt{vmm\_alloc\_pages} 内部想注册 VMA（跟踪新分配的虚拟地址区间），
VMA 的 \texttt{add} 又要 \texttt{kmalloc} 分配红黑树节点——
循环依赖！

解决方案：红黑树后端使用\textbf{静态节点池}（编译期 \texttt{static} 数组），
预分配固定数量的节点（如 256 个），完全绕开 \texttt{kmalloc}。
代价是 VMA 上限固定，但内核启动阶段区域数量很少，256 足够。
\end{whybox}

静态节点池的实现使用空闲链表，实现 $O(1)$ 的分配和释放：

\codefile{src/kernel/mm/vma\_rbtree.c}
\begin{lstlisting}[style=cstyle]
static rb_node_t  g_pool[RB_MAX_NODES];
static rb_node_t* g_free_head;   /* 空闲链表头 */

static void pool_init(void)
{
    g_free_head = &g_pool[0];
    for (unsigned i = 0; i < RB_MAX_NODES - 1; i++)
        g_pool[i].right = &g_pool[i + 1];
    g_pool[RB_MAX_NODES - 1].right = NULL;
}

static rb_node_t* pool_alloc(void)
{
    if (!g_free_head) return NULL;
    rb_node_t* n = g_free_head;
    g_free_head = n->right;
    n->color = RB_RED;   /* 新节点默认红色 */
    n->parent = n->left = n->right = NIL;
    return n;
}

static void pool_free(rb_node_t* n)
{
    n->right = g_free_head;
    g_free_head = n;
}
\end{lstlisting}

用 \texttt{right} 指针串联空闲节点，省掉额外的 \texttt{next} 字段。
分配时从头部取，释放时插回头部——均为 $O(1)$。

Maple Tree 后端同理，使用独立的节点池（\texttt{maple\_node\_t g\_node\_pool[64]}）
和 VMA 数据池（\texttt{vm\_area\_t g\_vma\_pool[256]}）。

% ============================================================================

